
\phantomsection
\section*{Введение}
\addcontentsline{toc}{section}{Введение}

\textbf{Актуальность темы.} Частные медицинские клиники ежедневно обрабатывают записи пациентов, расписание врачей, результаты приёмов, медицинские документы и обращения администраторов. При использовании бумажных журналов, отдельных таблиц или несвязанных программных сервисов возрастает вероятность ошибок при записи на приём, усложняется контроль загрузки врачей и затрудняется быстрый доступ к медицинской информации пациента. Поэтому актуальной является разработка веб-приложения, которое объединяет онлайн-запись, электронные медицинские карты и личные кабинеты пользователей в единой информационной системе.

\textbf{Степень разработанности проблемы.} На рынке существуют готовые медицинские информационные системы и пациентские сервисы, например Инфоклиника.RU, Medesk, ONDOC и Cliniccards. Они предоставляют инструменты для записи пациентов, ведения медицинских карт, работы с расписанием и коммуникации с пациентами~\cite{infoclinica,medesk,ondoc,cliniccards}. Однако такие решения являются закрытыми коммерческими продуктами, рассчитанными на внедрение в конкретной клинике и не предназначенными для изучения внутренней архитектуры в рамках учебного проекта. Это обосновывает необходимость разработки веб-приложения, в котором основные процессы частной клиники представлены в упрощённом, но целостном виде.

\textbf{Объект исследования} - процессы цифрового взаимодействия частной клиники с пациентами, врачами и администраторами.

\textbf{Предмет исследования} - программные средства, модели данных и пользовательские сценарии, применяемые при разработке веб-приложения для онлайн-записи, ведения медицинских карт и организации личных кабинетов.

\textbf{Целью данной работы} является проектирование веб-приложения для частной клиники, обеспечивающего онлайн-запись пациентов на приём, хранение медицинской информации, работу с расписанием врачей и разграничение доступа между пациентом, врачом и администратором.

Для достижения поставленной цели необходимо решить следующие задачи: выполнить обзор существующих программных продуктов, проанализировать инструменты веб-разработки и существующие технологические стеки, выбрать технологию, сформулировать требования к системе, спроектировать целевую аудиторию, функциональность, модель данных и макеты страниц.
